\exercisesection{诺特环的维数}

\begin{enumerate}
\def\labelenumi{\arabic{enumi})}
\tightlist
\item
  令 A 为环，U 为 Spec(A) 的开子集。
\end{enumerate}

\begin{enumerate}
\def\labelenumi{\alph{enumi})}
\item
  证明映射 \(\mathfrak{p} \mapsto \mathrm{V}(\mathfrak{p}) \cap \mathrm{U}\) 将 U 双射到 U 的闭不可约子集集合。
\item
  证明 dim(U) 是属于 U 的 A 的素理想链长度集合的上界。
\item
  令 \(\Phi\) 为属于 U 的 A 的极小素理想集合。证明有
\end{enumerate}

\[
\dim (\mathrm{U}) = \sup _ {\mathfrak {p} \in \Phi} \dim (\mathrm{V} (\mathfrak {p}) \cap \mathrm{U}).
\]

\begin{enumerate}
\def\labelenumi{\alph{enumi})}
\setcounter{enumi}{3}
\item
  以下假设 A 是局部环且 U ≠ Spec(A)。若 A 是诺特环，证明等式 \(\dim(\mathrm{U}) = \sup_{\mathfrak{p}\in\Phi}(\dim(\mathrm{A}/\mathfrak{p}) - 1)\)。
\item
  令 h 为满足 0 \textless{} h \textless{} dim(A) 的整数。令 \(E_{h}\) 为 \(p \in U\) 中满足 \(ht(p) = h\) 且 \(\dim(A/p) = \dim(A) - h\) 的素理想集合。若 A 是诺特环，证明 \(E_{h}\) 是无限集。证明当 A 不是诺特环时 \(E_{h}\) 可能是有限集（取 A 为高度为 2 的赋值环，令 U 为开集 \(\text{Spec}(A) - \{m_{A}\}\)，并取 h = 1）。
\end{enumerate}

\begin{enumerate}
\def\labelenumi{\arabic{enumi})}
\setcounter{enumi}{1}
\item
  令 A 为诺特整环。假设 A 不是域且 A 的分式域是有限生成 A-代数。则 A 是半局部环且维数为 1。（证明存在 A 中的非零非可逆元素 \(f\)，使得 \(A[f^{-1}]\) 是域，并化归为证明 \(\dim(A/(f)) = 0\)；否则 \(A/(f)\) 将有一个高度为 1 的素理想，其在 A 中的逆像 p 是高度为 2 的素理想；但维数为 2 的局部环 \(A_p\) 只能有有限多个素理想，这与第 3 节第 3 号命题 6 矛盾。）
\item
  \begin{enumerate}
  \def\labelenumii{\alph{enumii})}
  \tightlist
  \item
    令 A 为诺特局部环，令 p 为 A 的素理想。证明不等式 \(\dim_{\mathfrak{p}}(\mathbf{A}) \geqslant \mathrm{ht}(\mathfrak{p}) + \dim(\mathbf{A}/\mathfrak{p}) - 1\)。
  \end{enumerate}
\end{enumerate}

\begin{enumerate}
\def\labelenumi{\alph{enumi})}
\setcounter{enumi}{1}
\tightlist
\item
  证明存在维数为 3 的诺特局部环 R，具有素理想 q，使得 ht(q) = 1 且 dim(R/q) = 1。（使用习题 16，第~83页，取 \(\mathbf{R} = \mathbf{A}[\mathbf{U}]_{\mathbf{q}}\)，并令 q = pR。）令 \(f \in \mathfrak{m}_{\mathbb{R}} - \mathfrak{q}\)；证明 \(\dim(\mathbf{R}_{f}) = 2\)（使用习题 1，第~86页）；因此有 \(\dim_{\mathfrak{q}}(\mathbf{R}) = 2\)，而 \(\operatorname{ht}(\mathfrak{q}) + \dim(\mathbf{R}/\mathfrak{q}) - 1 = 1.\)
\end{enumerate}

\begin{enumerate}
\def\labelenumi{\arabic{enumi})}
\setcounter{enumi}{3}
\item
  令 \(k\) 为域，令 \(n\) 为整数 \(\geqslant 1\)。对于每个整数 \(r \in [0, n]\)，令 \(M_r\) 为 \(k(X_1, ..., X_n)\) 的向量子空间，它由形如 \(X_1^{\alpha_1}\ldots X_r^{\alpha_r}\) 的单项式生成，其中 \(\alpha_1, ..., \alpha_r\) 属于 \(Z\) 且 \(\alpha_r > 0\)。置 \(a_r = M_r + \cdots + M_n\)，其中 \(0 \leqslant r \leqslant n + 1\)。则 \(a_0\) 是 \(k(X_1, ..., X_n)\) 的子环，而 \(a_r\) 是 \(a_0\) 的素理想。置 \(A = (a_0)_{a_1}\)，并令 \(p_r = (a_r)_{a_1}\)，其中 \(1 \leqslant r \leqslant n + 1\)。证明局部环 \(A\) 的素理想只有 \(p_r\)，并且有 \(ht(p_r) = n + 1 - r\) 以及 \(dim(A) = n\)。此外，\(p_1\) 是主理想。
\item
  令 \(\rho: \mathbf{A} \to \mathbf{B}\) 为完备诺特局部环的局部同态，且满足 \([\kappa(\mathbf{B}): \kappa(\mathbf{A})] < \infty\)。置 \(n = \dim(\mathbf{B}/\mathfrak{m}_{\mathbf{A}}\mathbf{B})\)；证明存在局部 A-同态 \(u: \mathbf{A}[[\mathbf{T}_1, ..., \mathbf{T}_n]] \to \mathbf{B}\)，使 B 成为 \(\mathbf{A}[[\mathbf{T}_1, ..., \mathbf{T}_n]]\)-有限代数。（将 \(\mathbf{B}/\mathfrak{m}_{\mathbf{A}}\mathbf{B}\) 的一个极大正则序列提升到 B 中，并使用 III, § 3, 习题 18。）
\item
  令 \(k\) 为域，A 为形式幂级数环 \(k[[X_1, X_2]]\)。考察 A 的理想 \(\mathfrak{n}\)，它包含 A 的极大理想 \(\mathfrak{m}\) 的某个幂；以及由 \((X_1 - X_2)\mathfrak{n}\) 生成的理想 I。令 M 为 A/I-模。证明序列 \((X_1)\) 对 M 是正则的，但不是对 M 完全正则的。
\item
  令 A 为维数为 3 的局部整诺特非串联环（习题 16，第~83页），p ⊂ A 为满足 ht(p) = 1、dim(A/p) = 1 的素理想。证明存在 x ∈ p，使得 p 在包含 Ax 的素理想中是极小的。推出 V(x) 有一个维数为 1 的不可约分支。
\end{enumerate}

